\section{Related work}
\label{sec:related}

\subsection{Primitive Importance Score Prediction}
\label{sec:related_method}
Existing approaches for estimating Gaussian primitive importance can be broadly categorized into attribute-based, rendering-based, and learning-based methods.

\textbf{Attribute-based methods} adopt simple heuristics such as Gaussian volume and opacity, which were used in 3DGS~\cite{3dgs} and ADC-GS~\cite{adc} to detect redundant points. These methods are computationally efficient and easy to integrate, but they fail to reliably capture a Gaussian’s actual contribution to rendering quality, especially when multiple Gaussians overlap or interact in complex ways.  

\textbf{Rendering-based methods} derive importance by measuring each Gaussian’s contribution to rendered images, either through visibility or via reconstruction error sensitivity.
LightGaussian~\cite{lightgaussian} computes importance as the 2D projected area multiplied by absolute opacity and further refined by the 3D Gaussian volume. 
PRoGS~\cite{progs} and EAGLES~\cite{eagles} replace absolute opacity with blending opacity obtained during rasterization, which better reflects accumulated visibility but lacks volume refinement. 
MesonGS~\cite{mesongs} and HGSC~\cite{HGSC} combine blending opacity with volume, effectively integrating the strengths of both strategies for more stable scoring.
These strategies provide a closer approximation to perceptual importance but come with notable limitations. 
They rely on custom Gaussian rasterizers, and their computation time increase proportionally with the number of images used for score calculation. 
In addition, their performance is sensitive to the number and distribution of views, as shown in the Supplementary Material (Sec.~\ref{sup:render_views}).
Gradient-based variants estimate significance from reconstruction error sensitivity. 
C3DGS~\cite{c3dgs} uses backward gradients, selecting the maximum SH coefficient channel for color and the maximum eigenvalue of the rotation or scaling matrix for geometry, while PUP-3DGS~\cite{pup} extends this idea by computing perturbation sensitivity from the Hessian of reconstruction errors. 
While gradient signals can effectively indicate importance in well-optimized regions, they often produce unstable scores in poorly reconstructed areas, where large gradients arise from fitting errors rather than true contribution.

\textbf{Learning-based approaches} incorporate joint optimization of a per-Gaussian mask into the reconstruction pipeline, enabling data-driven selection of effective primitives \cite{compact,hac,hacpp,contextgs,maskgaussian}.
The mask implicitly encodes point significance, but these methods are tightly coupled with the reconstruction process, require lengthy optimization, and lack reusability. 
Moreover, once the Gaussian set changes (e.g., after pruning or editing~\cite{gaussianeditor, grouping,vc_edit}), the learned importance distribution becomes invalid, limiting generalization and transferability.

Motivated by these limitations, we propose RAP, a lightweight and rendering-free framework for efficient and generalizable Gaussian importance estimation.

\subsection{Applications of Primitive Importance Score}
\label{sec:app}
The concept of Gaussian primitive importance has been extensively explored in 3D reconstruction, where importance-guided strategies are employed to eliminate redundancy while preserving scene fidelity.
Prune-and-refine approaches \cite{lightgaussian,pup,hybridgs} typically remove low-importance primitives before refining the remaining set, while LP-3DGS \cite{lp3dgs} further learns to determine a favorable pruning ratio.

Alternatively, Mini-Splatting~\cite{minisplatting} selects representative Gaussians through sampling instead of pruning, mitigating potential artifacts from aggressive removal.
More recently, TAMING-3DGS~\cite{taming3dgs} utilizes importance scores to drive targeted densification, enhancing reconstruction efficiency under hardware and resource constraints.

Gaussian primitive importance has also become a key factor in compression and coding frameworks. 
Threshold-based pruning, adopted in MesonGS~\cite{mesongs} and EAGLES~\cite{eagles}, discards low-importance Gaussians prior to encoding to reduce data size. 
Beyond simple thresholding, C3DGS~\cite{c3dgs} employs sensitivity-aware clustering to group Gaussians according to their visual contribution, thereby preserving perceptually important structures while improving compression efficiency. 
Another direction focuses on learnable masking, with methods \cite{hac,hacpp,contextgs,maskgaussian} jointly optimizing mask values during reconstruction to guide the selection of effective primitives.
This idea, first introduced in Compact-3DGS~\cite{compact}, has  been extended in compression-oriented frameworks to enable adaptive importance modeling.

